\section{Discussion}
\label{sec:discussion}

\subsection{The Token-Direction Program: Partial Success, Fundamental Limitation}

Our results show that a shallow linear program over token directions can explain a meaningful fraction of MLP output variance---up to 54\% at layer~11 with $k=100$ directions, and 70\% with $k=500$.
At first glance, this seems to support the hypothesis that MLPs operate as simple programs in token space.

However, the random-direction control experiment (\tabref{tab:baseline}) reveals that this success is not specific to token directions.
Random orthogonal directions achieve comparable or superior $\Rsq$ values across most layers.
The implication is clear: the token-direction program works because the MLP is partially linear, not because the vocabulary provides a natural basis for describing MLP computation.
This distinction matters for interpretability: one cannot claim that the MLP ``thinks in tokens'' any more than it ``thinks in random directions.''

\subsection{The U-Shaped Linearity Profile}

The most striking structural finding is the U-shaped linearity profile across layers (\tabref{tab:linearity}).
Early layers (0--2) perform nearly linear transformations ($\Rsq > 0.81$).
Middle layers (3--8) are strongly nonlinear ($\Rsq$ drops to 0.10).
Late layers (10--11) partially linearize again ($\Rsq \approx 0.67$--$0.71$).

This pattern aligns with prior observations that early layers handle syntactic processing and late layers prepare final predictions~\citep{geva2021transformer,meng2022locating}, but adds a quantitative characterization.
The high nonlinearity of middle layers is consistent with these layers performing the most complex computation---disentangling superposed features~\citep{elhage2022toy}, composing multi-hop reasoning, and transforming representations in ways that require the full expressivity of the GELU nonlinearity.

\subsection{Distributed Computation}

The sparsity analysis (\tabref{tab:sparsity}) reveals that 50--66\% of MLP neurons are active at any given position, and the top-100 neurons explain only 18--23\% of the total activation.
This contradicts a simple ``key-value lookup'' interpretation where a handful of neurons fire for each input.
Instead, MLP computation is highly distributed, with thousands of neurons contributing small updates that collectively determine the output.
This distributed nature helps explain why simple programs fail in middle layers: the computation cannot be reduced to a few dominant directions.

\subsection{Implications for Mechanistic Interpretability}

Our findings have several implications for the field.

{\bf The ``MLP as token program'' metaphor has limited scope.}
It works reasonably for late layers preparing final predictions but breaks down for middle layers where the most interesting computation occurs.
Researchers using token-space projections to interpret MLP behavior should be aware that this lens captures at most half the picture, and that the captured portion reflects linearity rather than token-specific organization.

{\bf Random baselines are essential.}
Without the random-direction control, one might conclude that token directions are a meaningful basis for MLP computation.
We recommend that future interpretability studies include random-direction baselines when claiming that a particular basis (token directions, SAE features, or otherwise) is privileged.

{\bf Middle layers need nonlinear descriptions.}
The field should look beyond linear analyses---toward sparse autoencoders~\citep{cunningham2023sparse}, nonlinear probes, or other methods---to describe middle-layer MLP computation.
Linear methods, regardless of basis, capture less than 25\% of the variance in these layers.

\subsection{Limitations}
\label{sec:limitations}

\para{Model scope.}
We analyze only \gpttwo (12 layers, 85M parameters).
Larger models may exhibit different linearity profiles, and the relative importance of token directions may change with scale.

\para{Token direction definition.}
We define token directions via the unembedding matrix $\Wu$.
Alternative definitions---the embedding matrix $\mW_E$, SAE-derived features, or directions found by probing---might yield different results.

\para{Linear programs only.}
We test only linear (Ridge regression) programs.
Decision trees, polynomial features, or small neural networks might capture substantially more of the MLP behavior, especially in middle layers.

\para{Aggregate analysis.}
We analyze MLP behavior aggregated across diverse inputs.
The MLP might implement a simple program \emph{per input type} (e.g., factual recall prompts) while appearing complex in aggregate.
Input-conditional analysis is a promising direction.

\para{Top-$k$ selection.}
Using globally top-$k$ most variable token dimensions may miss input-specific relevant directions.
Per-input direction selection could improve results.
